% ============================================================================
% Introduction générale
% ============================================================================
\chapter*{Introduction générale}
\addcontentsline{toc}{chapter}{Introduction générale}
\markboth{Introduction générale}{}

Modifier une ligne de production réelle coûte cher~: chaque essai immobilise des machines, consomme de la matière et fait courir un risque sur la production. L'une des promesses de l'Industrie~4.0 est de déplacer une partie de ces essais vers un modèle virtuel fidèle de l'installation, le \textbf{jumeau numérique}, alimenté par des données et capable de répondre à la question \enquote{que se passerait-il si~?} avant toute intervention physique.

Ce mini-projet, troisième sujet de la liste proposée par Prof.~Maria Zemzami~\cite{zemzami2026}, a pour objectif de \textbf{modéliser une ligne de production simplifiée sous forme de jumeau numérique, permettant de simuler son comportement et de tester des scénarios d'optimisation}. Trois livrables sont attendus~: un modèle de simulation du système physique, une interface synchronisée avec des données simulées, et un scénario d'optimisation testé et documenté.

Nous avons réalisé ces trois livrables sous la forme d'une application complète, développée à six et versionnée dans un dépôt Git~: une simulation Python de la ligne, un backend temps réel qui la diffuse par WebSocket et en conserve l'historique, une scène \ac{3D} réalisée avec Unity et exportée en WebGL, un tableau de bord de supervision, ainsi qu'un module de comparaison de scénarios et une suite de tests automatisés qui valident la cohérence du modèle. Les résultats d'optimisation présentés dans ce rapport ont été mesurés avec la scène Unity elle-même, connectée au backend.

Le rapport suit le plan du cahier des charges. Le chapitre~\ref{chap:contexte} situe le projet dans l'Industrie~4.0 et en fixe le périmètre et les objectifs. Le chapitre~\ref{chap:specs} détaille les spécifications fonctionnelles et techniques. Le chapitre~\ref{chap:conception} présente la conception (architecture, modèle de simulation, protocole d'échange, répartition des rôles), et le chapitre~\ref{chap:implementation} sa réalisation. Le chapitre~\ref{chap:optimisation} documente le scénario d'optimisation et ses résultats chiffrés. Le chapitre~\ref{chap:organisation} traite des aspects organisationnels et humains, et le chapitre~\ref{chap:conclusion} dresse le bilan, les limites assumées et les perspectives.
